\chapter{Conclusion and Future Work}

\section{Conclusion}
This thesis set out to address a concrete and growing problem in digital entertainment: the fragmentation of discovery across two historically separate domains---film and gaming---and the information overload that fragmentation produces. The result of the work is \textbf{Questro}, a unified, socially-enriched, cross-domain recommendation platform that treats a user's cinematic and gaming tastes as a single coherent profile rather than two disconnected histories.

Across the preceding chapters we moved deliberately from problem to artifact. Chapter~1 framed the paradox of choice and motivated a cross-domain solution; Chapter~2 established the development methodology; Chapter~3 surveyed the field and identified the precise gap Questro fills; Chapter~4 formalized the system through requirement analysis and UML modelling; Chapter~5 specified the decoupled, microservice-oriented architecture; and Chapter~6 documented the concrete implementation in detail. Together these chapters demonstrate a complete engineering trajectory rather than an isolated prototype.

The technical contributions are anchored in two analytical centerpieces. The first is the \textbf{CL-EPIDTN} recommender, a dual-tower contrastive Transformer that embeds users and items into a shared space and is trained with a multi-objective criterion combining the main InfoNCE loss, a self-supervised contrastive loss over augmented views, and an auxiliary modality-matching loss. Trained on MovieLens-32M ratings and Steam play sessions---two corpora with \emph{zero} user overlap---the model performs a genuine cross-domain, cold-start semantic alignment, and at inference it is hybridized with content and user-based collaborative signals through reciprocal rank fusion. As reported in Table~\ref{tab:clepidtn_measured}, it ranks the single held-out item inside the top fifty of roughly $138{,}000$ candidates for nearly one user in five under a strict full-catalog, leave-one-out protocol. The second centerpiece is the \textbf{Retrieval-Augmented Generation chatbot}, which couples a memory-mapped FAISS index, parental-control filtering enforced at the vector boundary, personalized reranking, and Google Gemini grounding; across its 115-query evaluation suite (Table~\ref{tab:rag_eval}) it completed every query without error, mixed both domains in $68.7\%$ of responses, and surfaced the expected ground-truth item two-thirds of the time.

Surrounding these AI components, the platform delivers a production-grade ASP.NET Core backend built on Clean Architecture, the Specification Pattern, and an explicit-mapping persistence layer; a stateless dual-token security flow that is resistant to cross-site scripting token theft; deterministic, guardian-managed parental controls; and a responsive React~19 frontend whose silent token-refresh makes session management invisible to the user. The system therefore satisfies the objectives set out in Chapter~1: unified cross-domain discovery, activity tracking, child safety, and a social, community-driven experience, realized on a secure and performant technical foundation.

\section{Future Work}
Questro is a complete and functional system, but several avenues would meaningfully extend its data, intelligence, and reach. The most impactful of these directly attacks the central limitation of the current recommender; the remainder strengthen the retrieval pipeline and the broader platform.

\subsection{Enriching the Model with Physical-Media Reviews}
\label{sec:future_physical_reviews}
The single most valuable extension concerns the training data itself. The current corpus is drawn entirely from \textit{digital-first} sources---Steam play sessions and MovieLens ratings---which omits the large population of consumers who buy and review \textbf{physical copies} of games and films. A natural next step is to harvest ratings and textual reviews of boxed and physical editions from online retail and marketplace storefronts (for example, large e-commerce catalogues and second-hand marketplaces), normalize them against Questro's existing item catalogue, and fold them into the training corpus as additional interaction signals.

This enrichment is attractive for three reasons. First, it substantially increases the \textbf{volume and coverage} of interaction data, including pre-streaming-era titles and audiences that digital-only sources under-represent, which directly mitigates the data-sparsity and cold-start pressures discussed throughout this thesis \cite{ref:cold_start}. Second, and most importantly, retail reviewers frequently purchase \emph{both} films and games from the same storefront, which introduces the genuine \textbf{cross-domain user overlap} that the present MovieLens/Steam pairing structurally lacks; learning from users who demonstrably consume both media types would let CL-EPIDTN ground its cross-domain alignment in real shared behaviour rather than inferring it purely from content semantics \cite{ref:cross_domain_survey}. Third, the free-text reviews are themselves a rich signal: a sentiment-analysis stage could convert review prose into the same normalized $[-1, 1]$ interaction weights the model already consumes, turning qualitative opinion into quantitative training targets \cite{taha2024}.

Realizing this pipeline raises practical challenges that future work must address: respecting the terms of service and legal constraints of each data source, performing reliable entity resolution to deduplicate physical editions against existing TMDB/RAWG catalogue entries, filtering promotional and spam reviews, and de-biasing for the differing rating distributions of retail versus enthusiast platforms. Handled carefully, however, this single addition would expand both the size and the cross-domain richness of the training set more than any architectural change to the model.

\subsection{Retrieval-Augmented Generation Enhancements}
The RAG chatbot offers several clear directions for improvement:
\begin{itemize}
    \item \textbf{Conversational memory.} The pipeline is currently single-turn: each query is answered independently. Maintaining per-session dialogue state would let the assistant handle follow-up and refinement queries (``something lighter than that'', ``but for two players'') and progressively narrow recommendations across a conversation.
    \item \textbf{Hybrid retrieval.} Dense semantic search excels at meaning but can miss exact-title and rare-keyword matches. Fusing the existing dense FAISS retrieval with a sparse lexical retriever such as BM25 \cite{robertson2009bm25} would improve recall for precise queries while preserving semantic generalization.
    \item \textbf{Cross-encoder reranking.} Beyond the personalized reranker, a lightweight cross-encoder that jointly attends to the query and each candidate \cite{nogueira2019passage} could refine the ordering of the retrieved pool before personalization, improving the relevance of the candidate set the language model ultimately sees.
    \item \textbf{Streaming and latency reduction.} The measured end-to-end latency (Table~\ref{tab:rag_eval}) is dominated by the external language-model round-trip. Streaming tokens to the client as they are generated would sharply reduce \emph{perceived} latency, and a smaller or self-hosted grounding model would reduce absolute latency for high-traffic scenarios.
    \item \textbf{Response caching and feedback learning.} Caching embeddings and responses for popular queries would cut both latency and cost, while a lightweight thumbs-up/thumbs-down feedback channel would supply implicit labels for continuously improving retrieval and reranking quality.
\end{itemize}

\subsection{Recommender System Enhancements}
Several extensions would deepen the recommender beyond its current offline-trained form:
\begin{itemize}
    \item \textbf{Online and incremental learning.} The model is presently retrained offline and its catalogue baked into static artifacts. An incremental-update path would let newly observed Questro interactions refresh embeddings without a full retrain, keeping recommendations current as the live user base grows.
    \item \textbf{Native Questro cross-domain histories.} Once real users accumulate genuine cross-domain interaction sequences on the platform, the model can be fine-tuned on \emph{in-domain} data, complementing the public-corpus pre-training with behaviour observed on Questro itself.
    \item \textbf{Multimodal item representations.} Incorporating poster images, screenshots, and trailers through pre-trained vision encoders would enrich the item tower with visual semantics, complementing the current text-and-metadata content channel.
    \item \textbf{Online evaluation.} Offline ranking metrics tell only part of the story. Introducing controlled A/B testing and online engagement metrics (click-through, dwell time, conversion) would validate offline gains against real user satisfaction and support principled diversity/serendipity tuning.
\end{itemize}

\subsection{Platform and Infrastructure Enhancements}
Finally, the wider system would benefit from operational and feature investments:
\begin{itemize}
    \item \textbf{Scalability and deployment.} Containerizing the microservices and orchestrating them (e.g., with Docker and Kubernetes) behind autoscaling, together with a distributed cache (such as Redis) for sessions and one-time passwords, would support the concurrent-user targets defined in the non-functional requirements.
    \item \textbf{Observability.} End-to-end telemetry, structured logging, and distributed tracing across the backend, recommender, and RAG services would make performance regressions and failures diagnosable in production.
    \item \textbf{Security hardening.} Adopting cross-property validation through FluentValidation, request rate-limiting, and periodic penetration testing would extend the existing defense-in-depth posture.
    \item \textbf{Reach and engagement.} Native mobile clients, real-time social notifications, internationalization, and accessibility improvements would broaden Questro's audience and deepen the community engagement that distinguishes it from purely algorithmic recommenders.
\end{itemize}

\section{Closing Remarks}
Questro demonstrates that the two largest pillars of interactive entertainment---film and gaming---can be served from a single, coherent, and safety-aware recommendation platform, and that a contrastively-trained cross-domain model can be made to work end-to-end despite the absence of shared users between its public training corpora. The future directions outlined above, led by the enrichment of the training data with physical-media reviews, chart a clear path from this working system toward a richer, faster, and more deeply personalized successor.
